\documentclass[pdflatex,sn-apa]{sn-jnl}

\usepackage{graphicx}%
\usepackage{multirow}%
\usepackage{amsmath,amssymb,amsfonts}%
\usepackage{amsthm}%
\usepackage{mathrsfs}%
\usepackage[title]{appendix}%
\usepackage{xcolor}%
\usepackage{textcomp}%
\usepackage{manyfoot}%
\usepackage{booktabs}%
\usepackage{algorithm}%
\usepackage{algorithmicx}%
\usepackage{algpseudocode}%

\theoremstyle{thmstyleone}%
\newtheorem{theorem}{Theorem}%
\newtheorem{proposition}[theorem]{Proposition}%
\newtheorem{lemma}[theorem]{Lemma}%
\newtheorem{corollary}[theorem]{Corollary}%

\theoremstyle{thmstyletwo}%
\newtheorem{example}{Example}%
\newtheorem{remark}{Remark}%

\theoremstyle{thmstylethree}%
\newtheorem{definition}{Definition}%
\newtheorem{assumption}{Assumption}%

\newcommand{\indep}{\perp\!\!\!\perp}
\newcommand{\Ecal}{\mathcal{E}}
\newcommand{\Fcal}{\mathcal{F}}
\newcommand{\Scal}{\mathcal{S}}
\newcommand{\Gstar}{\mathcal{G}^{\star}}
\newcommand{\pa}{\mathrm{pa}}
\newcommand{\de}{\mathrm{de}}
\newcommand{\Prob}{\mathbb{P}}
\newcommand{\Exp}{\mathbb{E}}
\newcommand{\Real}{\mathbb{R}}

\raggedbottom

\begin{document}

\title[Certificate-Only Causal Discovery]{Certificate-Only Causal Discovery:
Anytime-Valid Adjacency and Orientation}

\author*[1]{\fnm{Quang-Vinh} \sur{Dang}}\email{vinh.dq4@buv.edu.vn}

\affil*[1]{\orgname{British University Vietnam}, \orgaddress{\city{Hung Yen},
\country{Vietnam}}}

\abstract{Constraint-based causal discovery removes an edge when a
conditional-independence test \emph{fails to reject}, converting absence of
evidence into evidence of absence. This is why finite-sample correctness for
such methods requires faithfulness, and it is why their output cannot
distinguish a feature the data support from one about which nothing was
learned. We develop a discovery algorithm built on a single rule: assert a
feature of the graph only by \emph{rejecting} a null that is true whenever
that feature is absent, and never assert anything by failing to reject. Two
certificates implement the rule. For adjacency, the statement ``some
conditioning set separates the pair'' is a union null, and the minimum of the
per-set e-processes is an e-process for it; crossing a threshold certifies an
edge under the Markov condition alone, with no faithfulness assumption. For
orientation, the two directions of a pair induce different joint densities
under a linear non-Gaussian model, so a sequential universal-inference
e-process against the reversed factorisation certifies an arrowhead, and
abstains automatically under Gaussian noise, where the direction is not
identified. Both certificates are anytime-valid, so the estimate may be
monitored continuously and the run stopped by any data-dependent rule. The
resulting algorithm, CERT-CD, grows from the empty graph, reports undecided
pairs as undecided, and bounds by $\alpha$ the probability that any certified
edge or arrowhead is \emph{ever} wrong, uniformly over time. On simulated
instances that violate $\delta$-strong faithfulness, the adjacency
certificates stay within budget while methods that delete on non-rejection
lose a true edge in most runs; the orientation certificates, whose null
additionally requires a correctly specified functional model, exceed theirs on
those same instances, and we identify the mechanism. On premise-satisfying
instances $96\%$ of certified edges are oriented, against a
Markov-equivalence-class ceiling of $33\%$. On the observational
protein-signalling data of Sachs and colleagues the two halves of the output
diverge instructively: every certified edge lies in the published consensus
skeleton, while every certified arrowhead is reversed relative to it---a
disagreement we trace to the linear non-Gaussian model rather than to the
inferential machinery, since DirectLiNGAM reverses the same arcs.}

\keywords{causal discovery, e-values, anytime-valid inference, safe testing,
family-wise error rate, LiNGAM}

\pacs[MSC Classification]{62H22, 62L12, 62F03}

\maketitle

\section{Introduction}\label{sec:intro}

A causal discovery algorithm returns a graph, and a practitioner reading that
graph would like to know which of its features the data support. Causal
graphical models have become a standard language for that question
\citep{pearl2009causality,peters2017elements}, and the algorithms that
estimate them from observational data are correspondingly consequential. Methods in
the constraint-based tradition---PC, FCI and their descendants
\citep{spirtes2000causation,spirtes1995causal,colombo2012learning}---answer this question
poorly, and the reason is structural rather than incidental. They begin from
the complete graph and delete an edge whenever a conditional-independence test
fails to reject. Failure to reject is not evidence of independence. It is
compatible with independence, with a dependence too weak for the test to see
at the current sample size, and with a test that had little power against the
alternative in question. The deletion happens regardless, and the resulting
absence of an edge is reported with the same authority as the presence of one.

Two consequences follow, and both are well known in their separate
literatures. First, finite-sample guarantees for these methods require
faithfulness or a quantitative strengthening of it, because only such an
assumption converts ``the test did not reject'' into ``the variables are
independent''. \citet{robins2003uniform} showed that no uniformly consistent
estimator of causal structure exists without such a strengthening, and
\citet{uhler2013geometry} showed that the set of parameter values violating
$\lambda$-strong faithfulness has substantial measure even in modest graphs.
The assumption is doing work the data cannot do. Second, the output of these
methods is two-valued where the evidence is three-valued: a missing edge means
either ``the data indicate no edge'' or ``nothing was established about this
pair'', and the graph does not say which.

This paper takes the opposite discipline. We adopt a single rule and follow it
to its consequences:

\begin{quote}
\emph{Assert a feature of the graph only by rejecting a null hypothesis that
is true whenever that feature is absent. Never assert anything by failing to
reject.}
\end{quote}

We call this the \emph{certificate-only} rule, and an assertion licensed by
such a rejection a \emph{certificate}. The rule is restrictive, and its
consequences are not cosmetic. An algorithm obeying it cannot begin from the
complete graph, because it has no licence to delete; it must begin from the
empty graph and add. It cannot report an edge as absent, only as undecided.
And the quantity it controls is the probability of asserting something false,
not the probability of missing something true. Power becomes a descriptive
property of a run rather than part of the guarantee.

Making the rule operational requires two certificates, and their statistical
situations are quite different.

\paragraph{Adjacency.} For a pair $(i,j)$, the null that is true whenever the
edge is absent is composite: \emph{some} conditioning set separates them. This
is a union over conditioning sets, and the minimum of the per-set e-processes
is an e-process for a union null (Section~\ref{sec:adjacency}). Two
consequences matter more than the construction itself. Validity uses no
faithfulness assumption---if $i$ and $j$ are non-adjacent in the true graph
then $\pa(i)$ or $\pa(j)$ separates them, so the null holds and Ville's
inequality applies---and the multiplicity correction runs over pairs rather
than over (pair, conditioning set) triples, because one hypothesis is tested
per pair.

\paragraph{Orientation.} Orientation error is not uncontrolled in the
literature. PC-p \citep{strobl2019estimating} attaches edge-specific p-values
to unshielded colliders and Meek-rule orientations and controls their false
discovery rate; \citet{shimizu2008use} test the direction of causation between
two variables using non-normality; \citet{strieder2023confidence} construct
confidence statements for causal effects under structure uncertainty. What has
not been available is orientation error control that is simultaneously
finite-sample, time-uniform, family-wise over a whole graph, and not confined
to the Markov equivalence class. Constraint-based orientation tests are
confined to that class by construction, so an edge lying in no v-structure is
unorientable in principle by any of them. The functional-model methods that
escape the class---LiNGAM and its descendants
\citep{shimizu2006linear,shimizu2011directlingam,hyvarinen2013pairwise}---choose
a direction by comparing scores, and the comparison itself carries no error
statement. Under a linear non-Gaussian model the two orientations of a pair
induce different joint densities and exactly one is correct, so ``the
direction is $j \to i$'' is a genuine null hypothesis; a sequential
universal-inference e-process against it certifies the arrowhead $i \to j$
(Section~\ref{sec:direction}). Under Gaussian noise the two factorisations are
observationally indistinguishable, the process does not grow, and no
certificate is issued---the procedure abstains exactly where the direction is
unidentified, without being told to.

\paragraph{Why anytime-valid.} Both certificates are built as e-processes
rather than as fixed-sample tests, and the choice is not decorative. A
discovery run on streaming or accumulating data invites the analyst to look at
the current graph, and looking is a form of optional stopping. A fixed-sample
guarantee is void the moment the analyst stops on the basis of what they saw.
Section~\ref{sec:exp-calibration} quantifies the gap: at nominal
$\alpha = 0.05$, a Fisher-$z$ test inspected after every observation rejects a
true null in $34.6\%$ of runs, against $5.6\%$ when it is evaluated once. An
e-process pays for time-uniformity in power, and we report that price rather
than eliding it.

\paragraph{Contributions.} The paper makes four.

\begin{enumerate}
\item An e-process for the composite \emph{separability} null, obtained as a
  minimum over conditioning sets, giving time-uniform certification of
  adjacency under the Markov condition and a bound on separator size, with no
  faithfulness assumption (Section~\ref{sec:adjacency}). We give the exact
  per-set primitive for the linear-Gaussian case in closed form, together with
  the confidence sequence obtained by inverting it.

\item A sequential certificate for causal \emph{orientation}
  (Section~\ref{sec:direction}) that is valid at any stopping time, is not
  confined to the Markov equivalence class, and abstains automatically under
  Gaussianity. We identify a requirement the construction imposes on the
  fitting procedure---the null-model fit must genuinely attain its
  supremum---which is easy to violate and which silently inflates the
  e-process when violated (Section~\ref{sec:supremum}).

\item CERT-CD, an algorithm combining both certificates under one budget
  (Section~\ref{sec:algorithm}). Its output is a three-valued graph in which
  every asserted edge and arrowhead carries a certificate and undecided pairs
  are explicit, with a whole-graph guarantee that is proved rather than
  assumed.

\item An empirical study (Section~\ref{sec:experiments}) that stratifies
  instances by whether the premise of the competing guarantee holds, rather
  than pooling them; measures the assumptions of our own theorem instead of
  asserting them; reports the regimes in which the certificates are void
  (Section~\ref{sec:exp-violations}); and applies the method to the
  observational protein-signalling data of \citet{sachs2005causal}.
\end{enumerate}

The study is not uniformly favourable to the method, and the two places where
it is not are the more informative ones. On instances violating
$\delta$-strong faithfulness the orientation certificates exceed their
family-wise budget (Section~\ref{sec:exp-main}), because the conditioning
block they freeze is built from already-certified neighbours and is therefore
incomplete exactly when certification is slow. On the protein-signalling data
every certified arrowhead is reversed relative to the published consensus
(Section~\ref{sec:exp-sachs}), and DirectLiNGAM reverses the same arcs, which
locates the failure in the shared functional model rather than in the
inference. Both are instances of one gap---the direction null assumes a model,
and a certificate is only as good as its null---and we have tried to give them
the prominence a guarantee-bearing method owes them.

We are explicit throughout about which components are standard.
Section~\ref{sec:related} gives that accounting in detail, including three
claims of novelty we made in earlier drafts and then withdrew after reading
the prior art.

\paragraph{Notation and outline.} Section~\ref{sec:background} fixes notation
and reviews e-processes and the graphical setting. Section~\ref{sec:principle}
states the certificate-only rule formally and derives its structural
consequences. Sections~\ref{sec:adjacency} and~\ref{sec:direction} develop the
two certificates. Section~\ref{sec:algorithm} assembles CERT-CD and proves the
whole-graph guarantee; Section~\ref{sec:latent} extends it to latent
confounders. Section~\ref{sec:related} positions the work.
Section~\ref{sec:experiments} reports experiments and
Section~\ref{sec:discussion} discusses limitations. Proofs and derivations
omitted from the main text are in the appendices.

\section{Background}\label{sec:background}

\subsection{E-variables, e-processes and Ville's inequality}
\label{sec:eproc}

Let $(\Omega, \Fcal, (\Fcal_t)_{t \ge 0}, \Prob)$ be a filtered probability
space and let $H_0$ be a set of distributions on it.

\begin{definition}[E-variable]\label{def:evar}
A non-negative random variable $E$ is an \emph{e-variable} for $H_0$ if
$\Exp_P[E] \le 1$ for every $P \in H_0$.
\end{definition}

\begin{definition}[E-process]\label{def:eproc}
A non-negative, adapted process $(E_t)_{t \ge 0}$ is an \emph{e-process} for
$H_0$ if $\Exp_P[E_\tau] \le 1$ for every $P \in H_0$ and every stopping time
$\tau$ (finite or not, with $E_\infty := \liminf_t E_t$).
\end{definition}

Every non-negative supermartingale starting at one is an e-process, by
optional stopping \citep{shafer2019game}; the converse fails, and the distinction matters for
composite nulls, where an e-process need not be a supermartingale under any
single null distribution \citep{ramdas2023gametheoretic}. Both of our
constructions are e-processes that are dominated pointwise by a
$P$-supermartingale for each $P \in H_0$, with the dominating process
depending on $P$. That is the standard route to an e-process for a composite
null, and it is the route both proofs below take.

The operational content of Definition~\ref{def:eproc} is Ville's inequality
\citep{ville1939etude}: for an e-process $(E_t)$ and any $\alpha \in (0,1)$,
\begin{equation}\label{eq:ville}
\Prob\bigl( \exists\, t \ge 0 : E_t \ge 1/\alpha \bigr) \;\le\; \alpha
\qquad \text{for every } P \in H_0 .
\end{equation}
The quantifier is inside the probability. This is what licenses continuous
monitoring: an analyst may inspect the evidence after every observation and
stop whenever they wish, and the probability that the process \emph{ever}
crosses the threshold under the null is still at most $\alpha$. A p-value
offers no analogue. A fixed-sample p-value controls the rejection probability
at one pre-specified sample size, and the guarantee is destroyed by looking;
the phenomenon is familiar from sequential clinical trials and from online
experimentation \citep{johari2022always}.

E-values also compose in ways p-values do not, and we use two such
composition rules. If $E^{(1)}, \dots, E^{(m)}$ are e-variables for
$H_0^{(1)}, \dots, H_0^{(m)}$ then $\min_k E^{(k)}$ is an e-variable for
$\bigcup_k H_0^{(k)}$, and any convex combination $\sum_k w_k E^{(k)}$ with
$w_k \ge 0$, $\sum_k w_k = 1$ is an e-variable for $\bigcap_k H_0^{(k)}$,
under \emph{arbitrary} dependence between the components
\citep{vovk2021evalues}. The first rule is what makes the adjacency
certificate possible; the second is what makes the multiplicity discussion of
Section~\ref{sec:multiplicity} non-trivial.

Recent work has developed e-processes for a wide range of settings---bounded
means \citep{waudbysmith2024estimating}, linear models
\citep{lindon2026anytime}, stratified counts \citep{turner2023safe}---and has
characterised admissible constructions \citep{ramdas2023gametheoretic}. Two
general recipes recur below. The first is the \emph{method of mixtures}: under
a group-invariant null, integrating the likelihood ratio against the right-Haar
measure on the nuisance group yields an exact test martingale
\citep{berger1998bayes,hendriksen2021optional,grunwald2024safe}.
Section~\ref{sec:primitive} instantiates it for partial correlation. The second
is \emph{universal inference} \citep{wasserman2020universal}: for a null model
$\{D_\theta : \theta \in \Theta_0\}$ and any alternative fitted on past data,
the ratio of a sequentially-fitted alternative likelihood to the maximised null
likelihood is an e-process, with no regularity conditions on $\Theta_0$ beyond
the existence of the supremum. Section~\ref{sec:direction} instantiates it for
causal direction. Universal inference is known to be conservative
\citep{tse2022note}, and Section~\ref{sec:exp-orientation} measures how
conservative it is here.

\subsection{Causal discovery setting}\label{sec:setting}

Let $\Gstar$ be a directed acyclic graph (DAG) with vertex set
$V = \{1,\dots,d\}$ identified with random variables $X_1,\dots,X_d$, and let
$P$ be the joint law of $X = (X_1,\dots,X_d)$. We write $\pa(i)$ for the
parents of $i$ in $\Gstar$ and $\de(i)$ for its descendants. Throughout,
observations $o_1, o_2, \dots$ are drawn i.i.d.\ from $P$ and arrive in
batches; $\Fcal_t$ is the $\sigma$-field generated by the first $t$
observations.

\begin{assumption}[Markov condition]\label{ass:markov}
$P$ is Markov to $\Gstar$: whenever $S$ d-separates $i$ from $j$ in $\Gstar$,
$X_i \indep X_j \mid X_S$ under $P$.
\end{assumption}

Assumption~\ref{ass:markov} is the direction of the correspondence that
follows from the causal structure itself: it says that the graph's separations
are reflected in the distribution. Its converse---faithfulness, that every
conditional independence in $P$ corresponds to a d-separation in
$\Gstar$---does not follow from the causal structure, and is the assumption we
do not make for validity.

Write $\Scal_k(i,j) = \{ S \subseteq V \setminus \{i,j\} : |S| \le k \}$ for
the family of admissible conditioning sets. This family is determined by $d$
and $k$ and is fixed before any data are seen; nothing about it adapts to the
sample. We say the pair $(i,j)$ is \emph{separable within $\Scal_k$} under $P$
if $X_i \indep X_j \mid X_S$ for some $S \in \Scal_k(i,j)$.

\begin{assumption}[Separator size]\label{ass:sep}
Every non-adjacent pair of $\Gstar$ has a d-separating set of size at most
$k$.
\end{assumption}

Assumption~\ref{ass:sep} is structural rather than distributional: it
constrains $\Gstar$, not $P$, and it can be guaranteed by choosing $k$ large
enough. It holds whenever $k$ is at least the maximum in-degree of $\Gstar$,
since for non-adjacent $i,j$ either $j \notin \de(i)$, in which case $\pa(i)$
d-separates them, or $i \notin \de(j)$, in which case $\pa(j)$ does. Taking
$k = d-2$ makes it vacuous at the cost of an exponential family $\Scal_k$; the
trade-off is computational, not inferential, and we return to it in
Section~\ref{sec:complexity}. Section~\ref{sec:exp-premise} measures how often
Assumption~\ref{ass:sep} holds at the $k$ we use, and compares it with how
often $\delta$-strong faithfulness holds on the same instances.

\subsection{Faithfulness and its quantitative strengthenings}
\label{sec:faithfulness}

Since the argument of this paper is about what faithfulness is doing, it is
worth being precise about the versions in play.

\emph{Faithfulness} holds if every conditional independence in $P$ is implied
by d-separation in $\Gstar$; \citet{meek1995strong} showed that it fails only
on a Lebesgue-null set of parameters for a fixed graph, which is the result
usually cited in its defence and which \citet{uhler2013geometry} showed to be
misleading in finite samples. \emph{Adjacency-faithfulness}
\citep{ramsey2006adjacency} is the weaker requirement that adjacent variables
are dependent given every conditioning set. \emph{$\delta$-strong
faithfulness} \citep{uhler2013geometry,kalisch2007estimating} is the
quantitative strengthening that every partial correlation not forced to zero
by d-separation has absolute value at least $\delta$; it is what turns
consistency into a finite-sample statement, and it is the premise of the
comparators in Section~\ref{sec:exp-main}.

The asymmetry we exploit is this. Faithfulness in any of its forms is an
assumption about which \emph{dependences} exist. A procedure that asserts an
edge by rejecting an independence null needs it only for \emph{power}: without
it, a genuinely adjacent pair whose dependence is cancelled will fail to be
certified, and the pair stays undecided. A procedure that asserts a
\emph{non}-edge by failing to reject needs it for \emph{validity}: without it,
a cancelled dependence produces a confidently deleted edge. The same
assumption sits on opposite sides of the guarantee depending on which
direction of the inference the algorithm runs. \citet{zhang2008detection} and
\citet{ramsey2006adjacency} respond to this by detecting or tolerating
unfaithfulness within a delete-on-non-rejection architecture; we respond by
changing the architecture.
